\documentclass[11pt]{article}
\usepackage[T1]{fontenc}
\usepackage{textcomp}
\usepackage[margin=0.75in]{geometry}
\usepackage{amsmath,amssymb,amsthm}
\usepackage{array,booktabs}
\usepackage{hyperref}
\setlength{\parindent}{0pt}
\newtheorem{theorem}{Theorem}
\theoremstyle{definition}
\newtheorem{definition}{Definition}
\title{Flow Proof Artifact: Subgroup-inv-closed}
\author{Generated by \texttt{flow doc proof}}
\date{Flow Proof Book}
\begin{document}
\maketitle
Subgroup inverse closure follows from group laws.\\[0.5em]
\textbf{Source.} dummit-foote — *Abstract Algebra*, §2.2\\[0.5em]
\bigskip
\noindent{\large \textbf{Derived fact 1.} \textit{If g in H then inv(g) in H when H is a subgroup} \hfill Subgroup \cdot inverse \cdot subgroup inverses stay in the subgroup}\par
\smallskip\noindent\textit{Coordinate: Subgroup \cdot inverse \cdot subgroup inverses stay in the subgroup}\par
\smallskip\noindent\textit{Source: dummit-foote}\par
\smallskip\noindent\textit{Built on: closed under inverses, for inverse on Subgroup, left inverse recovers the identity, for inverse on Group}\par
\medskip
\noindent\textbf{Goal.} If g in H then inv(g) in H when H is a subgroup\par
\noindent$\displaystyle \forall H \in Subgroup \forall g \in Group\quad inv(g) in H$\par
\medskip
\noindent
\begin{tabular}{@{} >{\raggedright\arraybackslash}p{0.06\textwidth} >{\raggedright\arraybackslash}p{0.47\textwidth} >{\raggedleft\arraybackslash}p{0.41\textwidth} @{}}
\textbf{\#} & \textbf{Proof} & \textbf{Mathematics} \\
\midrule
\textbf{1.} & We prove that subgroup inverses stay in the subgroup for inverse on Subgroup. & \\[0.45em]
\textbf{2.} & We split into exhaustive cases — the claim must hold in each one. & \\[0.45em]
\textbf{3.} & Case 1 (see step 2): suppose g in H. & \\[0.45em]
\textbf{4.} & We invoke the definitional clause governing inverse on Subgroup: closed under inverses, for inverse on Subgroup (instantiated for H, g). & \\[0.45em]
\textbf{5.} & We invoke the definitional clause governing inverse on Group: left inverse recovers the identity, for inverse on Group (instantiated for g). & \\[0.45em]
\textbf{6.} & From step 3, step 4, and step 5, this implies inv(g) in H. Together with the other cases (step 3), the goal is discharged. Hence proven. & $\displaystyle inv(g) in H$ \\[0.45em]
\bottomrule
\end{tabular}
\label{thm:Subgroup-inverse-subgroup_inverses_stay_in_the_subgroup}
\par\medskip\hrule
\end{document}
